% Industry resume, Grover Aruquipa
% Compile with pdfLaTeX (also works on Overleaf). One column on purpose:
% applicant tracking systems read single-column text far more reliably than sidebars.
\documentclass[10pt,letterpaper]{article}
\usepackage[T1]{fontenc}
\usepackage[utf8]{inputenc}
\usepackage[scaled=0.95]{helvet}
\renewcommand{\familydefault}{\sfdefault}
\usepackage[top=1.1cm,bottom=1.0cm,left=1.5cm,right=1.5cm]{geometry}
\usepackage{enumitem}
\usepackage{titlesec}
\usepackage{xcolor}
\usepackage{microtype}
\usepackage[hidelinks]{hyperref}

\definecolor{ink}{HTML}{111111}
\definecolor{accent}{HTML}{0D2B61}
\definecolor{muted}{HTML}{555555}
\color{ink}
\pagestyle{empty}
\setlength{\parindent}{0pt}
\setlength{\parskip}{0pt}

\titleformat{\section}{\normalsize\bfseries\color{accent}}{}{0em}{\MakeUppercase}[\vspace{-6pt}{\color{accent}\rule{\linewidth}{0.6pt}}]
\titlespacing*{\section}{0pt}{7pt}{3pt}
\setlist[itemize]{leftmargin=1.1em,itemsep=0.5pt,topsep=1pt,parsep=0pt,label=\textbullet}

% \role{Title}{Organization}{Location}{Dates}
\newcommand{\role}[4]{\textbf{#1}, #2 \hfill {\color{muted}#3 \textbar\ #4}\par}

\begin{document}

%================ HEADER ================
{\LARGE\bfseries Grover Aruquipa}\par\vspace{2pt}
{\large Robotics Researcher \textbar\ Robotic Grasping, Manipulation and Machine Learning on Real Hardware}\par\vspace{3pt}
{\small Qu\'ebec, QC, Canada \ \textbullet\ \
\href{mailto:grover.aruquipa.9@gmail.com}{grover.aruquipa.9@gmail.com} \ \textbullet\ \
\href{https://www.linkedin.com/in/grover-aruquipa-996917126}{linkedin.com/in/grover-aruquipa-996917126} \ \textbullet\ \
\href{https://groveraruquipa.github.io}{groveraruquipa.github.io}}\par

%================ SUMMARY ================
\section{Summary}
Builder-researcher finishing a PhD in Mechanical Engineering (Universit\'e Laval, defense January 2027) on robotic
grasping and manipulation. I derive the models on paper and then build the robot: three multi-DoF manipulators with
integrated grippers taken from theory and CAD to 500\,Hz real-time C++ control and experiments, with 3 journal papers,
4 more under review, and a review of robotic grasping in ASME JMR. Earlier industrial R\&D in vision and machine
learning (Hitachi ABB Power Grids, Sencon) means I know both the data side and the robot that has to execute.

%================ SKILLS ================
\section{Technical Skills}
\begin{itemize}
\item \textbf{Grasping and manipulation:} grasp analysis and grasp-quality indices, robots with integrated grippers,
unlimited in-hand rotation of grasped objects, contact and force transmission, redundancy resolution.
\item \textbf{Machine learning and perception:} PyTorch, TensorFlow, OpenCV; computer vision for industrial anomaly
detection; learning on multi-sensor industrial data (vibration and circuit-breaker sensors).
\item \textbf{Robot modeling and control:} rigid-body kinematics and dynamics, Lie groups and screw theory, Jacobians,
singularity and workspace analysis; multi-axis real-time control at 500\,Hz, trajectory generation and tracking.
\item \textbf{Software and simulation:} C++ (real-time control), Python, MATLAB and Simulink, ROS and ROS~2, MoveIt,
Pinocchio, MuJoCo, Simscape, Git, Linux; SolidWorks and prototype fabrication.
\end{itemize}

%================ EXPERIENCE ================
\section{Experience}
\role{Doctoral Researcher}{Robotics Laboratory, Universit\'e Laval}{Qu\'ebec, Canada}{2023 -- present}
\begin{itemize}
\item Designed, built and controlled three parallel manipulators (4, 7 and 8 actuators) with the gripper driven through
the mechanism, allowing unlimited rotation of the grasped object.
\item Wrote the real-time multi-motor control at 500\,Hz; validated a 4-DoF coaxial wrist-gripper over ten continuous
tool revolutions with sub-degree tracking while it held the object.
\item Derived closed-form inverse kinematics, Jacobians and singularity conditions, and a redundancy strategy that keeps
unlimited rotation free of collisions and singularities.
\item Currently building a Delta robot with an integrated gripper that estimates grasp force from motor currents, aimed
at handling fragile objects at speed; developing a ROS~2 and MoveIt~2 stack for a 9-DoF parallel robot.
\item Led two graduate students from concept to working prototype on an 8-actuator (6+2)-DoF robot (ICRA 2027
submission); reviewer for ASME J.\ Mechanisms and Robotics.
\end{itemize}

\vspace{3pt}
\role{Visiting Researcher}{Institute of Robotics, Johannes Kepler University Linz}{Linz, Austria}{Apr. -- May 2026}
\begin{itemize}
\item Developed singularity and grasp-quality indices for robots that grasp with an articulated platform, with
Univ.-Prof.~A.~M\"uller (Mitacs Globalink Research Award).
\end{itemize}

\vspace{3pt}
\role{Research Intern}{FEMTO-ST (CNRS)}{Besan\c{c}on, France}{2022 -- 2023}
\begin{itemize}
\item Modeled a redundant dexterous parallel robot for micro-manipulation and wrote its real-time C++ control software.
\end{itemize}

\vspace{3pt}
\textbf{Industry R\&D in vision and machine learning} \hfill {\color{muted}2020 -- 2021}\par
\begin{itemize}
\item \textbf{Graduate Research Intern}, Hitachi ABB Power Grids R\&D (Zurich, 2020): machine learning on sensor data from
high-voltage circuit breakers; invited after my bachelor's thesis.
\item \textbf{Computer Vision Engineer}, Sencon Srl (2021): anomaly detection on industrial production lines. Also data
scientist for KTH Royal Institute of Technology (2021).
\end{itemize}

%================ EDUCATION ================
\section{Education}
\textbf{PhD, Mechanical Engineering}, Universit\'e Laval \hfill {\color{muted}Qu\'ebec, Canada \textbar\ 2023 -- Jan.\ 2027 (expected)}\par
{\small Kinematically redundant parallel robots with unlimited rotations and integrated grasping. Advisors: C.~Gosselin, P.~Lambert.}\par\vspace{2pt}
\textbf{MSc, Smart Integrated Systems}, Universit\'e Bourgogne--Franche-Comt\'e \hfill {\color{muted}Besan\c{c}on, France \textbar\ 2021 -- 2023}\par\vspace{2pt}
\textbf{BEng, Mechatronics Engineering}, Bolivian Catholic University \hfill {\color{muted}La Paz, Bolivia \textbar\ 2014 -- 2019}\par
{\small Thesis: reconfigurable Delta robot for product picking.}

%================ AWARDS ================
\section{Awards and Languages}
\begin{itemize}
\item 1st Place, Graduate Division, ASME Student Mechanisms and Robotics Design Competition (2025), team lead.
\item Finalist, MassRobotics competition, Boston, with one of my robot prototypes, team lead.
\item Mitacs Globalink Research Award (2026); 1st Place, Student Paper Competition, Argentine Robotics Conference (2019).
\item Languages: Spanish (native), English (full professional), French (intermediate).
\end{itemize}

\newpage
%================ PUBLICATIONS ================
\section{Publications}
{\small First author on all papers.}\par\vspace{2pt}
\textbf{Journal articles}
\begin{itemize}
\item \textbf{G.~Aruquipa}, P.~Lambert, R.~Dahmouche, C.~Gosselin, ``Parallel robots for grasping tasks: a review,''
\emph{ASME Journal of Mechanisms and Robotics}, vol.~18, no.~12, art.~120801, 2026. doi:10.1115/1.4072612
\item \textbf{G.~Aruquipa}, V.~Muralidharan, P.~Lambert, C.~Gosselin, ``Analysis and design of a novel compact
four-degree-of-freedom parallel robot with unlimited rotation,'' \emph{ASME Journal of Mechanisms and Robotics},
vol.~18, no.~6, art.~061001, 2026.
\item \textbf{G.~Aruquipa}, P.~Lambert, C.~Gosselin, ``Analysis and design of a novel (6+1)-DoF parallel robot with
unlimited rotation and grasping capabilities,'' \emph{Mechanism and Machine Theory}, vol.~224, art.~106438, 2026.
\end{itemize}
\textbf{Conference papers}
\begin{itemize}
\item \textbf{G.~Aruquipa}, P.~Lambert, C.~Gosselin, ``Kinematic analysis and design of a novel 9-DoF parallel robot
with grasping capabilities,'' \emph{CCToMM Symposium on Mechanisms, Machines and Mechatronics}, 2025, pp.~117--127.
\item \textbf{G.~Aruquipa}, F.~Diaz, ``An IoT architecture based on the control of a bio-inspired manufacturing system
for the detection of anomalies with vibration sensors,'' \emph{Procedia Computer Science}, vol.~200, pp.~438--450, 2022.
\end{itemize}
\textbf{Under review}
\begin{itemize}
\item \textbf{G.~Aruquipa}, P.~Lambert, C.~Gosselin, ``A kinematically redundant (6+2)-DoF parallel robot with three
orthogonal unlimited rotations and coaxial actuation,'' submitted to \emph{IEEE ICRA 2027}.
\item \textbf{G.~Aruquipa}, P.~Lambert, C.~Gosselin, ``A compact coaxial parallel wrist-gripper with unlimited
tool-axis rotation,'' submitted to \emph{Robotics and Autonomous Systems}.
\item \textbf{G.~Aruquipa}, P.~Lambert, C.~Gosselin, ``An analytical framework for collision-free and singularity-free
unlimited rotation in spatial parallel mechanisms,'' submitted to \emph{Mechanism and Machine Theory}.
\item \textbf{G.~Aruquipa}, A.~Ayoub, R.~Dahmouche, ``Architecture, kinematics, and workspace analysis of a redundant
8RSS parallel robot with two unlimited rotations,'' in revision, \emph{Mechanism and Machine Theory}.
\end{itemize}


%================ REFERENCES ================
\section{References}
\begin{itemize}
\item \textbf{Prof.~Cl\'ement Gosselin}, Full Professor, Department of Mechanical Engineering, Universit\'e Laval, Qu\'ebec,
Canada. PhD supervisor. \href{mailto:clement.gosselin@gmc.ulaval.ca}{clement.gosselin@gmc.ulaval.ca}
\item \textbf{Prof.~Patrice Lambert}, Assistant Professor, Department of Mechanical Engineering, Universit\'e Laval,
Qu\'ebec, Canada. PhD co-supervisor. \href{mailto:patrice.lambert@gmc.ulaval.ca}{patrice.lambert@gmc.ulaval.ca}
\item \textbf{Univ.-Prof.~Dr.-Ing.~Andreas M\"uller}, Head of the Institute of Robotics, Johannes Kepler University Linz,
Austria. Host during my research stay (2026). \href{mailto:a.mueller@jku.at}{a.mueller@jku.at}
\end{itemize}

\end{document}
